\documentclass[11pt,a4paper]{article}

\usepackage[T1]{fontenc}
\usepackage[utf8]{inputenc}
\usepackage{lmodern}
\usepackage[a4paper,margin=25mm]{geometry}
\usepackage{amsmath,amssymb,amsthm,mathtools}
\usepackage{booktabs,array}
\usepackage{float}
\usepackage{microtype}
\usepackage{xcolor}
\usepackage{enumitem}
\usepackage{url}
\usepackage[colorlinks=true,linkcolor=blue!55!black,citecolor=blue!55!black,
            urlcolor=blue!55!black]{hyperref}

\newtheorem{theorem}{Theorem}
\newtheorem{lemma}[theorem]{Lemma}
\newtheorem{remark}[theorem]{Remark}
\newcommand{\F}{\mathbb{F}}
\newcommand{\Hthree}{\mathsf{H}_3}
\newcommand{\Hash}{\mathsf{H}}
\newcommand{\Sign}{\mathsf{Sign}}
\newcommand{\Verify}{\mathsf{Verify}}
\newcommand{\pk}{\mathsf{pk}}
\newcommand{\sk}{\mathsf{sk}}
\newcommand{\concat}{\mathbin\Vert}

\title{\bfseries Omitted Fiat--Shamir Binding Enables Universal Forgery\\
\large A Public-Key-Only Break of All 14 Submitted Chinith Parameter Sets}
\author{Martin Feussner\\
\small Selmer Center, University of Bergen\\
\small \texttt{martin.feussner@uib.no}}
\date{25 September 2026}

\begin{document}
\maketitle

\begin{center}
\fbox{\parbox{0.92\linewidth}{\small\textbf{Provenance and review status.}
The attack was independently found by an AI system, OpenAI Codex using
Daybreak Blue at maximum reasoning effort, without a human first proposing
this specific attack.  The derivation, implementation, experiments, and
initial manuscript were also produced by the AI.  At the date above, no human
cryptanalyst has independently verified the argument, experiments, or novelty
assessment.  This is a preliminary disclosure for human review and
reproduction.}}
\end{center}
\vspace{0.7em}

\begin{abstract}
Chinith is a family of VOLE-in-the-head signatures submitted to the Chinese
Next-generation Commercial Cryptographic Algorithms competition.  Its
QuickSilver proof reduces satisfaction of the public one-way-function
relation to a compressed constant term $\widetilde a_0$.  The specification
binds that value into the final Fiat--Shamir challenge together with
$\widetilde a_1,\ldots,\widetilde a_{d-1}$.  Both submitted implementation
trees instead hash only $\widetilde a_1,\ldots,\widetilde a_{d-1}$.  The
signer computes $\widetilde a_0$ and discards it.  The verifier reconstructs
$\widetilde a_0$ from the public relation and proof transcript, then also
discards it.  Consequently, verification never tests whether the committed
witness satisfies the relation represented by the public key.

This yields a universal forgery requiring no signature queries and no secret
information.  Given any public key and chosen message, an attacker selects an
arbitrary witness and arbitrary seed material, executes the public prover
operations, derives the defective final challenge, and emits an ordinary
signature.  Every remaining commitment, consistency, grinding, and message
binding check is honestly satisfied; the sole missing check is the one that
would reject the false witness.

We implemented this attack against the submitted portable code.  The
untouched reference verifiers accepted public-key-only forgeries for every one
of the 14 SM4th, uBlockith, and Vistrutith parameter sets.  Forging took
0.011--2.660 single-threaded CPU-seconds.  Complete processes, including key
generation solely to obtain a legitimate target public key and two verifier
calls, took 0.07--7.92 seconds and at most 55 MB of resident memory on a
3.2 GHz Intel Core i5-6500.  Changing one bit of the target message caused
every forged signature to be rejected.  A static audit finds the same omission
in all 14 reference and all 14 optimized implementations.  The written
specification contains the missing binding, so the result is a complete break
of the submitted implementations rather than an attack on the correctly
specified transcript.
\end{abstract}

\section{Introduction}

Chinith comprises the SM4th, uBlockith, and Vistrutith signature families.
Each proves knowledge of a block-cipher key or input satisfying a public
one-way-function relation.  The submission builds this proof from
VOLE-in-the-head commitments and a QuickSilver-style arithmetic proof, then
applies the Fiat--Shamir transform to obtain a non-interactive signature
\cite{Chinith,VOLEitH,QuickSilver,FiatShamir}.

The final Fiat--Shamir challenge has two jobs.  It selects the VOLE openings,
and it binds the compressed QuickSilver proof into the transcript.  For
degree-$d$ constraints, the prover obtains field elements
$\widetilde a_0,\ldots,\widetilde a_{d-1}$.  The verifier reconstructs
$\widetilde a_0$ from its own evaluation of the public relation.  It then
hashes that reconstructed value together with the transmitted higher
coefficients and checks that the result equals the opening challenge.  This
fixed-point check is what makes a false witness fail.

The submitted code omits exactly the reconstructed term.  On both sides, the
final hash begins at the serialized address of $\widetilde a_1$ and consumes
$(d-1)$ field elements.  Thus it hashes
$\widetilde a_1,\ldots,\widetilde a_{d-1}$ and never
$\widetilde a_0$.  This is consistent between signer and verifier, so known
answer tests generated by the same code do not reveal the error.

The omission removes the public-key relation from acceptance.  An attacker can
run the prover with the all-zero witness, or any other bit string of the
required length, while retaining the victim's public key as the public input
to the constraint system.  The VOLE layer correctly commits to that false
witness.  The constraint layer computes evidence that it is false, but that
evidence appears only in $\widetilde a_0$, which acceptance ignores.

Our contributions are:

\begin{itemize}[leftmargin=1.5em]
  \item We identify the omitted constant proof term in all 28 submitted
        implementation directories: 14 reference and 14 optimized variants.
  \item We give a public-key-only universal forgery.  It uses no signing
        oracle, no recovered key, no malformed input, and no rare event.
  \item We prove acceptance by separating the still-valid VOLE transcript
        checks from the discarded QuickSilver relation check.
  \item We reproduce the forgery end to end against the untouched reference
        verifier for all 14 parameter sets, including the degree-2 and
        degree-3 proof paths and all three block-cipher families.
  \item We describe a direct repair and regression tests that distinguish a
        corrected implementation from the submitted one.
\end{itemize}

\section{The relation check in Chinith}
\label{sec:scheme}

We summarize only the elements needed for the attack.  Let the public key be
$\pk=(x,y)$.  Depending on the variant, the claimed relation is either
\begin{equation}
                         y=E_k(x)
\label{eq:standard-owf}
\end{equation}
or the Even--Mansour relation
\begin{equation}
                         y=E_x(k)+k.
\label{eq:em-owf}
\end{equation}
Here and below addition is over the relevant characteristic-two space.  The
signer extends the secret $k$ into a bit witness $w\in\{0,1\}^{\ell}$ that
contains the internal values needed to check the selected block cipher.

\subsection{VOLE commitment and masked witness}

The prover obtains a VOLE correlation consisting of a bit vector $u$, masks
$v_i\in\F_{2^{\lambda_F}}$, verifier values $q_i$, and a global challenge
$\Delta\in\F_{2^{\lambda_F}}$ satisfying
\begin{equation}
                         q_i=u_i\Delta+v_i.
\label{eq:vole}
\end{equation}
It sends the masked witness
\begin{equation}
                              D=w+u.
\label{eq:masked-witness}
\end{equation}
After reconstruction, the verifier can therefore treat
$q_i+D_i\Delta$ as a commitment to $w_i$ with mask $v_i$.

The first two Fiat--Shamir challenges bind the message, public key, VOLE
commitment, consistency data, and $D$.  None of those steps requires $w$ to
satisfy Equation~\eqref{eq:standard-owf} or~\eqref{eq:em-owf}; they establish
only that the later proof refers to the same committed bit string.

\subsection{Compressed QuickSilver check}

Chinith evaluates the public constraint system on the VOLE-committed witness.
For a degree-$d$ relation, each compressed constraint has the polynomial form
\begin{equation}
  \rho(X)=a_0+a_1X+\cdots+a_{d-1}X^{d-1}+zX^d,
\label{eq:constraint-poly}
\end{equation}
where the leading coefficient $z$ vanishes when the constraint is satisfied.
A universal hash seeded by the second challenge batches all constraints and
masking terms into
\[
             \widetilde a_0,\widetilde a_1,\ldots,
             \widetilde a_{d-1}\in\F_{2^{\lambda_F}}.
\]
The verifier evaluates its committed view at
$\Delta=\mathsf{ToField}(\mathsf{chall}_3)$ and derives a value
$\widetilde q$.  In the submission's characteristic-two notation it
reconstructs
\begin{equation}
  \widetilde a_0^{\star}
    =\widetilde q+
      \sum_{i=1}^{d-1}\widetilde a_i\Delta^i.
\label{eq:a0-reconstruction}
\end{equation}
Using Equation~\eqref{eq:constraint-poly}, the reconstructed constant term can
be written as
\begin{equation}
  \widetilde a_0^{\star}
    = \widetilde a_0 + \widetilde z\,\Delta^d,
\label{eq:a0-discrepancy}
\end{equation}
where $\widetilde z$ is the compressed leading coefficient of the batched
constraint polynomial.  For a valid committed witness, $\widetilde z=0$ and
therefore $\widetilde a_0^{\star}=\widetilde a_0$.  For an invalid witness,
the discrepancy in Equation~\eqref{eq:a0-discrepancy} is exactly the value
that the final Fiat--Shamir binding is supposed to expose; equality fails
except with the intended soundness probability.

\subsection{The specified final transcript}

The specification derives the final challenge inside its grinding loop as
\begin{equation}
 \mathsf{chall}_3=Hthree(
    \mathsf{chall}_2\concat
    \widetilde a_0\concat\widetilde a_1\concat\cdots\concat
    \widetilde a_{d-1}\concat\mathsf{ctr}).
\label{eq:correct-hash}
\end{equation}
The signature transmits $\widetilde a_1,\ldots,\widetilde a_{d-1}$ but omits
$\widetilde a_0$, because the verifier reconstructs it with
Equation~\eqref{eq:a0-reconstruction}.  Correct verification substitutes
$\widetilde a_0^{\star}$ into Equation~\eqref{eq:correct-hash} and accepts only
when the resulting challenge equals the transmitted $\mathsf{chall}_3$.

This creates the necessary dependency cycle for soundness.  The final
challenge determines $\Delta$ and hence the verifier's reconstructed constant
term, while the constant term is itself included in the hash that must produce
the final challenge.  An honest proof satisfies the cycle.  A prover using a
false witness cannot simply derive the challenge before learning the value at
which its inconsistent constraint polynomial will be checked.

\section{The submitted omission}
\label{sec:bug}

The implementation allocates and computes $\widetilde a_0$ during signing.
It serializes only the higher terms, as the compact signature format requires.
When initializing the final hash, however, it supplies a pointer to the first
serialized term $\widetilde a_1$ and hashes exactly $d-1$ field elements.  Its
effective challenge is
\begin{equation}
 \mathsf{chall}_3=Hthree(
    \mathsf{chall}_2\concat
    \widetilde a_1\concat\cdots\concat
    \widetilde a_{d-1}\concat\mathsf{ctr}).
\label{eq:defective-hash}
\end{equation}
The locally computed $\widetilde a_0$ is freed without affecting any
serialized byte or challenge.

Verification makes the matching error.  It evaluates the block-cipher
constraints and computes $\widetilde a_0^{\star}$ using
Equation~\eqref{eq:a0-reconstruction}.  It then initializes the final hash at
the signature's $\widetilde a_1$ field and consumes the same $d-1$ elements.
The reconstructed value is freed without being compared or hashed.

The code pattern appears in every parameter directory.  It is shared across
the four standard SM4th variants, four SM4th-EM variants, two standard
uBlockith variants, two uBlockith-EM variants, and two Vistrutith variants.
The same signer and verifier pattern occurs in both the portable reference and
optimized implementation trees.  The degree does not protect any variant:
the degree-2 path hashes only $\widetilde a_1$, while the degree-3 paths hash
$\widetilde a_1\concat\widetilde a_2$.

Equation~\eqref{eq:defective-hash} contains no result of the verifier's public
relation evaluation.  The block-cipher constraint code still runs, but its
output has no path to the acceptance decision.

\subsection{Code-level evidence}
\label{sec:code-evidence}

The relevant source-level behavior can be summarized by the following
normalized pseudocode.  The notation is schematic: it records the data flow
observed in the submitted signer and verifier without depending on local
variable names or memory-layout details.

\begin{verbatim}
Specification:
    chall3 = H3(chall2 || a0 || a1 || ... || a[d-1] || ctr)

Submitted signer:
    (a0, a1, ..., a[d-1]) = OWF_Prove(...)
    chall3 = H3(chall2 || a1 || ... || a[d-1] || ctr)

Submitted verifier:
    a0_star = ReconstructRelationCheck(...)
    chall3_star = H3(chall2 || a1 || ... || a[d-1] || ctr)
    accept only checks chall3_star == chall3
\end{verbatim}

Thus the signer computes the constant relation-check term and the verifier
reconstructs it, but neither value is absorbed into the final challenge.  The
same omission is present in the portable reference and optimized trees.  An
independent source audit should record the exact file names, function names,
and line or commit locations for these three operations when reproducing the
result; this report does not invent source identifiers that have not yet been
independently checked.

\section{Universal forgery}
\label{sec:attack}

The attack uses the ordinary public algorithms.  The arbitrary value denoted
$k^{\star}$ below is used only as input material for deriving commitment
randomness; it is not asserted to be a preimage of the public key.

\paragraph{Input.}
A public key $\pk=(x,y)$ and a chosen target message $M$.

\paragraph{Forgery algorithm.}
\begin{enumerate}[leftmargin=1.7em]
  \item Choose any bit string $w^{\star}\in\{0,1\}^{\ell}$, any string
        $k^{\star}$ of the block-cipher key length, and any salt $\rho$.
        All-zero $w^{\star}$ and $k^{\star}$ suffice.
  \item Compute the message digest from $\pk$ and $M$.  Use $k^{\star}$, the
        digest, and $\rho$ in the submitted randomness derivation to obtain a
        root seed and IV.
  \item Run the ordinary VOLE commitment.  Derive $\mathsf{chall}_1$, compute
        the VOLE consistency response, and set
        $D=w^{\star}+u$.  Derive $\mathsf{chall}_2$ exactly as the signer does.
  \item Run the ordinary, publicly specified OWF prover routine on $w^{\star}$ and the victim public key.
        It returns
        $(\widetilde a_0,\ldots,\widetilde a_{d-1})$.  Retain the serialized
        values $\widetilde a_1,\ldots,\widetilde a_{d-1}$; the value
        $\widetilde a_0$ is irrelevant to the submitted verifier.
  \item Perform the normal counter search using
        Equation~\eqref{eq:defective-hash}.  Decode the resulting challenge
        and produce a valid BAVC opening, retrying the counter exactly as in
        ordinary signing.
  \item Return the normal signature encoding of the commitment corrections,
        consistency response, $D$, higher proof coefficients, opening,
        $\mathsf{chall}_3$, IV preimage, and counter.
\end{enumerate}

Every operation above is public and has essentially the cost of one signing
operation.  No transcript from a legitimate signer is used.

\begin{theorem}[Acceptance of the forgery]
For every honestly generated Chinith public key $\pk$, every target message
$M$, and every chosen witness $w^{\star}$ of the required format, whenever
the ordinary grinding-and-opening procedure succeeds, the submitted verifier
accepts the resulting signature independently of whether $w^{\star}$ satisfies
the public OWF relation.
\end{theorem}

\begin{proof}
The commitment and opening are generated honestly, so VOLE reconstruction
succeeds.  Both parties recompute the same first challenge because they use
the same public key, message digest, commitment, correction values, and IV.
Equation~\eqref{eq:masked-witness} and the valid VOLE correlation imply that
the verifier's consistency computation represents the same chosen witness
$w^{\star}$.  It therefore reconstructs the same second challenge.

The verifier next evaluates the public OWF constraints on its committed view
and obtains some $\widetilde a_0^{\star}$.  This value need not equal the
prover's $\widetilde a_0$ because $w^{\star}$ may be false.  Under the submitted
acceptance rule, neither value is used.  The only proof coefficients absorbed
by the final hash are the transmitted
$\widetilde a_1,\ldots,\widetilde a_{d-1}$.  The verifier hashes those same
bytes with the already matching $\mathsf{chall}_2$ and counter.  It therefore
recomputes exactly the attacker's $\mathsf{chall}_3$.  The grinding predicate
and BAVC opening were selected for that challenge, so all remaining checks
also pass.
\end{proof}

\begin{remark}
The target message remains bound to the forged signature through the earlier
challenges.  This does not impede forgery because the attacker chooses the
message before producing the transcript.  It does predict that changing the
message after forging should cause rejection, which our experiments confirm.
\end{remark}

\section{End-to-end evaluation}
\label{sec:experiments}

We evaluated the attack on the official Chinith submission archive with
SHA-256 digest
\begin{center}
\small\ttfamily
4d31ba3fe8718f59901b7efdf01b4179912ef8e04f58b4b6dbba84186f002283.
\end{center}
The experiments used the portable reference implementation, GCC 13.3.0, and
the submission's release-style optimization and assertion settings.  The
machine was a four-core Intel Core i5-6500 at 3.2 GHz with 16 GiB of RAM,
running 64-bit Linux.  Each forgery itself ran on one thread.

For each parameter set, we first generated a valid key pair solely to obtain a
legitimate target public key.  We then overwrote the secret key and the inputs
from which it was generated.  From that point onward the attack received only
the public key, a 37-byte target message, a chosen salt, an all-zero fake key,
and an all-zero witness of the parameter-prescribed length.  It called the
public prover operations and passed the result to the untouched submitted
verification API.  We separately evaluated the one-way function under the
all-zero fake key and confirmed that its output differed from the target
public-key output in every run.

For a negative control, we flipped the first bit of the message and invoked
the same verifier on the same forged signature.  In all 14 cases, verification
returned success for the target and failure for the altered message.  Table
\ref{tab:results} gives the serialized signature size, CPU time spent in the
forgery, whole-process elapsed time, and peak resident memory.  Whole-process
time includes target key generation, forgery, successful verification, and
the altered-message verification.

\begin{table}[H]
\centering
\caption{Public-key-only forgeries against all submitted reference variants.
Every target signature was accepted and every one-bit message alteration was
rejected.}
\label{tab:results}
\scriptsize
\setlength{\tabcolsep}{3.3pt}
\begin{tabular}{@{}lrrrrr@{}}
\toprule
Variant & $d$ & Signature & Forge CPU & Total wall & Peak RSS \\
        &     & (bytes)   & (seconds) & (seconds)  & (MiB) \\
\midrule
SM4th-d3-128f-loose       & 3 &   6,724 & 0.043 & 0.12 &  3.1 \\
SM4th-d3-128f-tight       & 3 &  11,864 & 0.067 & 0.15 &  3.8 \\
SM4th-d3-128s-loose       & 3 &   5,056 & 0.123 & 0.34 & 10.3 \\
SM4th-d3-128s-tight       & 3 &   9,176 & 0.145 & 0.78 & 15.1 \\
SM4th-EM-d2-128f-loose    & 2 &   4,932 & 0.011 & 0.07 &  2.5 \\
SM4th-EM-d2-128f-tight    & 2 &   9,450 & 0.012 & 0.07 &  3.1 \\
SM4th-EM-d2-128s-loose    & 2 &   3,818 & 0.071 & 0.50 & 10.6 \\
SM4th-EM-d2-128s-tight    & 2 &   7,556 & 0.081 & 0.71 & 11.6 \\
uBlockith-d3-256f         & 3 &  31,556 & 0.066 & 0.25 & 12.8 \\
uBlockith-d3-256s         & 3 &  24,144 & 0.305 & 1.63 & 16.6 \\
uBlockith-EM-d3-256f      & 3 &  25,028 & 0.048 & 0.31 &  3.7 \\
uBlockith-EM-d3-256s      & 3 &  19,056 & 0.260 & 2.26 & 16.0 \\
Vistrutith-d3-512f        & 3 & 106,788 & 0.552 & 2.84 &  8.5 \\
Vistrutith-d3-512s        & 3 &  83,004 & 2.660 & 7.92 & 54.6 \\
\bottomrule
\end{tabular}
\end{table}

The measurements show a stronger result than laptop-scale key recovery.  The
adversary needs zero signing queries, less than three CPU-seconds for the
largest forgery, and memory measured in tens of megabytes.  Signature sizes
are identical to ordinary signatures for their respective parameter sets.

As a diagnostic before constructing the public-key-only attack, we also used
the SM4th-d3-128s-loose prover with an honestly generated public key and
modified each of the 226 encoded witness-byte positions individually.  All
226 resulting proofs were accepted.  This experiment is diagnostic only: the
modified byte strings need not all correspond to semantically valid alternative
witnesses.  The all-zero-witness experiments in Table~\ref{tab:results} are
the stronger end-to-end result because they demonstrate acceptance for a
single explicit false witness whose one-way-function output was separately
checked not to match the target public key.

\section{Repair and validation}

The direct repair is to implement Equation~\eqref{eq:correct-hash}.  During
signing, the final-hash state must absorb the locally computed
$\widetilde a_0$ and then every serialized field element
$\widetilde a_1,\ldots,\widetilde a_{d-1}$ before absorbing the counter.
During verification, it must absorb the reconstructed
$\widetilde a_0^{\star}$ followed by the same serialized higher coefficients.
Changing only the starting pointer is insufficient because
$\widetilde a_0$ is stored separately from the signature buffer.

The repaired challenge changes every signature, so submitted known-answer
vectors will need regeneration.  Keys and the public signature layout need
not change: $\widetilde a_0$ remains reconstructible and therefore need not be
transmitted.

At minimum, validation should include the following negative tests:
\begin{itemize}[leftmargin=1.5em]
  \item generate the prover transcript with an arbitrary witness under an
        unrelated valid public key and require rejection;
  \item change every witness byte in turn, regenerate the proof, and require
        rejection except when the modified value independently satisfies the
        relation;
  \item compare the exact byte sequence absorbed into the signer's and
        verifier's final hashes and confirm that it begins with the respective
        $\widetilde a_0$ value;
  \item run cross-implementation known-answer tests generated from a corrected
        transcript definition rather than accepting agreement between two
        implementations that share the same helper.
\end{itemize}

Including $\widetilde a_0$ restores the intended QuickSilver relation check,
but the repaired implementation and its security claims should be reviewed in
full.  This report does not claim that the correctly specified construction
has a comparable forgery.

\section{Scope and disclosure status}

The result applies to the Chinith archive identified in
Section~\ref{sec:experiments}.  It covers every submitted parameter set and
both implementation trees at the source level; the end-to-end measurements
use the portable reference tree.  It does not depend on a crash, undefined
input, uninitialized memory, weak entropy, timing leakage, or a signing
oracle.  The signature is produced for an attacker-chosen message and is
accepted by the ordinary verifier for the victim public key.

The mathematical algorithms in the submitted specification explicitly place
$\widetilde a_0$ in the final hash.  Accordingly, our claim is a universal
forgery against the submitted software realizations and their generated test
vectors.  We do not claim a break of an implementation that follows the
specified transcript.

The attack was independently found by OpenAI Codex using Daybreak Blue at
maximum reasoning effort, without a human first proposing this specific
attack.  Human verification, independent reproduction, and novelty review
remain pending.

\section{Conclusion}

Chinith's submitted signer and verifier agree on a final Fiat--Shamir
transcript that excludes the one reconstructed field element carrying the
QuickSilver relation check.  That agreement makes ordinary signatures verify
and self-generated known-answer tests pass, while allowing any committed
witness to stand in for the secret key.  The resulting attack is a
zero-query, public-key-only universal forgery across all 14 submitted
parameter sets, with sub-second forging time for 13 variants and 2.660 seconds
for the largest measured variant.

\begin{thebibliography}{9}

\bibitem{Chinith}
Chinith submission team.
\newblock \emph{Chinith: Algorithm Specifications}.
\newblock NGCC submission package, 2026.

\bibitem{QuickSilver}
K.~Yang, P.~Sarkar, C.~Weng, and X.~Wang.
\newblock QuickSilver: Efficient and Affordable Zero-Knowledge Proofs for
  Circuits and Polynomials over Any Field.
\newblock In \emph{ACM CCS}, pages 2986--3001, 2021.

\bibitem{VOLEitH}
C.~Baum, L.~Braun, C.~D. de Saint Guilhem, M.~Kloo\ss, E.~Orsini,
  L.~Roy, and P.~Scholl.
\newblock Publicly Verifiable Zero-Knowledge and Post-Quantum Signatures from
  VOLE-in-the-Head.
\newblock In \emph{CRYPTO}, pages 581--615, 2023.

\bibitem{FiatShamir}
A.~Fiat and A.~Shamir.
\newblock How to Prove Yourself: Practical Solutions to Identification and
  Signature Problems.
\newblock In \emph{CRYPTO}, pages 186--194, 1986.

\bibitem{FAEST}
C.~Baum, W.~Beullens, L.~Braun, et al.
\newblock \emph{FAEST v2: Algorithm Specifications, Version 2.0}.
\newblock 2025.

\end{thebibliography}

\end{document}
